%% kapa: exam=1 round=21 year=2010 period=1 form=A subject=economics subject_ko=경제원론 seq=35 number=75
%% source: Q-Net 감정평가사 기출문제 (cst00309) artlSeq=5024181 file=extracted/1008/5024181/2010년 제21회 감정평가사 1차시험 기출문제/1교시.hwp
%% conversion: kordoc 4.24.51 -> tools/gen1.py (auto); review=proofread-sonnet; flags=none
\begin{question}{75}
다음의 경제모형에서 균형 국민소득과 균형 이자율은 각각 얼마인가? (단, Y는 국민소득, C는 소비, I는 투자, G는 정부지출, L은 화폐수요, M은 화폐공급, r은 이자율이다.)
\begin{framedtext}
$Y = C + I + G$ (생산물시장의 균형식) \par
$L = M$ (화폐시장의 균형식) \par
$C = 100 + 0.8Y$ \qquad $I = 50 - 2r$ \qquad $G = 150$ \par
$L = 0.2Y - 2r$ \qquad $M = 100$
\end{framedtext}
\par\smallskip
\hspace{1.2cm}\uline{\makebox[2.6cm]{균형 국민소득}}\hspace{1em}\uline{\makebox[2.2cm]{균형 이자율}}\par
\begin{choices}
  \item \makebox[2.6cm][c]{600}\hspace{1em}\makebox[2.2cm][c]{10}
  \item \makebox[2.6cm][c]{700}\hspace{1em}\makebox[2.2cm][c]{20}
  \item \makebox[2.6cm][c]{800}\hspace{1em}\makebox[2.2cm][c]{30}
  \item \makebox[2.6cm][c]{900}\hspace{1em}\makebox[2.2cm][c]{40}
  \item \makebox[2.6cm][c]{1,000}\hspace{1em}\makebox[2.2cm][c]{50}
\end{choices}
\end{question}
\answer{5}
